\section{Problem 2.}

A right triangle has a line drawn from the vertex of the right angle to the hypotenuse, making two more right angles --
this is called the perpendicular.  This perpendicular divides the triangle into two segments that have lengths of 9 and 16 respectively. 
Find the length of the perpendicular, and the two legs of the triangle as shown below.

\begin{figure}[h]
    \centering
    \includegraphics[width=0.25\linewidth]{p2.png}
\end{figure}

\begin{solution}
There are many ways to solving this and the easiest is probably the usage
of similair triangles in combination with some Pythaogrean theorem. Look at
how the altitude splits the right angle into two complementary angles. From this
relationship, we also know that the other smaller triangle will share the same angles
and therefore the same ratios hence similiar triangles. Using this, we can find the 
altitude length.

\begin{figure}[h]
    \centering
    \begin{minipage}[t]{0.45\textwidth}
        \vspace{0pt}
        Let $a$ be the altitude side $\overline{BE}$.
        \begin{align*}
            \frac{9}{a} &= \frac{a}{16} \\
            (9)(16) &= a^{2} \\
            a &= 12, \ \text{lengths are only positive.}
        \end{align*}
    \end{minipage}
    \begin{minipage}[t]{0.45\textwidth}
        \vspace{0pt}
        \centering
        \includegraphics[width=0.5\linewidth]{labeledp2.jpg}
        \caption{$\triangle ABC \sim \triangle AEB \sim \triangle BEC$}
    \end{minipage}
\end{figure}

Now we can use Pythaogrean Theorem on both sides.
\begin{align*}
    \overline{AB}^{2} &= 9^{2} + 12^{2} \\
    \overline{AB} &= 15 \\
    \overline{BC}^{2} &= 16^{2} + 12^{2} \\
    \overline{BC} &= 20
\end{align*}

In whole, these are all of our lengths. $\boxed{\overline{BE} = 12, \ \overline{AB} = 15, \ \overline{BC} = 20}$

\subsection{Extra Methods}

Just thought it was interesting with the variety of methods. My actual initial 
thought was to create a system of equation and spam Pythaogrean Theorem.

\begin{align*}
    \overline{BE}^{2}+9^{2} &= \overline{AB}^{2} \\
    \overline{BE}^{2}+16^{2} &= \overline{BC}^{2} \\
    \overline{AB}^{2}+\overline{BC}^{2} &= 25^{2}
\end{align*}

Solving this isn't actually that bad because we can substitue all of the squared 
varables with just another variable and then unsquare it when we are done.

\begin{align*}
    \begin{bmatrix}
            -1 & 0 & 1 \\
            0 & -1 & 1 \\
            1 & 1 & 0
    \end{bmatrix}
    \begin{bmatrix}
        \overline{AB}^{2} \\
        \overline{BC}^{2} \\
        \overline{BE}^{2}
    \end{bmatrix}
    &=
    \begin{bmatrix}
        -81 \\
        -256 \\
        625
    \end{bmatrix} \\
    \begin{bmatrix}
        \overline{AB}^{2} \\
        \overline{BC}^{2} \\
        \overline{BE}^{2}
    \end{bmatrix}
    &=
    \begin{bmatrix}
        225 \\
        400 \\
        114
    \end{bmatrix} \\
    \begin{bmatrix}
        \overline{AB} \\
        \overline{BC} \\
        \overline{BE}
    \end{bmatrix}
    &=
    \begin{bmatrix}
        15 \\
        20 \\
        12
    \end{bmatrix} 
\end{align*}

Another method quite literally is the exact usage of the Geometric mean theorem
in which we can solve for the altitude.
\begin{figure}[h]
    \centering
    \includegraphics[width=0.4\linewidth]{geothm.png}
    \caption{From Wikipedia $h=\sqrt{pq}$.}
\end{figure}

\begin{align*}
    h &= \sqrt{16 \cdot 9} \\
    \Longrightarrow h &= 12 = \overline{BE}
\end{align*}
From which, we can simply solve for both side legs in the same way as the original method.

\end{solution}